% --- Archon physics formalization source begin ---
% archon:physics
% archon:covers PhyXMiniProblems/problem_phyx_mini_0334.lean
% archon:source-report reports/phyx_mini/problem_phyx_mini_0334.source.json

\chapter{Physics problem phyx\_mini\_0334}
\label{ch:PhyXMiniProblems_problem_phyx_mini_0334}

\paragraph{Problem source.}
\#\# Physical scenario

A large tank of water has a hose connected to it. The tank is sealed at the top and has compressed air between the water surface and the top. When the water height \$h\$ has the value 3.50\textbackslash{},m, the absolute pressure \$p\$ of the compressed air is \$4.20 \textbackslash{}times 10\textasciicircum{}5\$\textbackslash{},Pa. Assume that the air above the water expands at constant temperature, and take the atmospheric pressure to be \$1.00 \textbackslash{}times 10\textasciicircum{}5\$\textbackslash{},Pa. As water flows out of the tank, \$h\$ decreases.

\#\# Question

At what value of \$h\$ does the flow stop?

\#\# Answer choices

- A: A: 1.74m

- B: B: 1.55m

- C: C: 1.64m

- D: D: 1.92m

\#\# Figure caption (auxiliary; use the image as primary evidence)

The image depicts a cylindrical container filled with a liquid. The container has a height of 4.00 meters. It features a side opening connected to a pipe, which is positioned 1.00 meter above the base of the container. 

In the diagram:

- The letter “p” is labeled at the top of the liquid, indicating pressure at the surface.
- The letter “h” is marked from the surface of the liquid down to the bottom of the container, signifying the height of the liquid.
- Arrows are used to denote vertical measurements from the top to the bottom of the cylinder and from the bottom to the opening for the pipe.
- Text labels provide the specific lengths 4.00 m for the container height and 1.00 m for the distance from the base to the pipe outlet.

\#\# Dataset metadata

Ideal Gases and Kinetic Theory; Multi-Formula Reasoning, Physical Model Grounding Reasoning

\paragraph{Recorded answer/context.}
A: A: 1.74m

\paragraph{Figure/image path.}
/root/proposal\_for\_physic/science-mango/phyx\_mini\_run/phyx\_data/test\_image/334.png

\paragraph{Formalization target.}
create a compiling Lean file with sorry bodies at `PhyXMiniProblems/problem\_phyx\_mini\_0334.lean`. The Lean declarations must preserve the physical quantities, dimensions or dimensional roles, figure labels, governing-law hypotheses, and final relation expressed by this problem.
Use Mathlib/Physlib names found through LeanExplore where available. If a physics API is missing, introduce faithful local abstractions rather than scalar placeholder aliases.

\begin{theorem}[Physics formalization target]
\label{thm:physics:phyx_mini_0334:target}
\lean{PhyXMiniProblems.ProblemPhyXMini0334.problem_phyx_mini_0334}
\uses{def:physics:phyx-mini-0334:phyxminiproblems-problemphyxmini0334-lengthquantity, def:physics:phyx-mini-0334:phyxminiproblems-problemphyxmini0334-lengthinmeters, def:physics:phyx-mini-0334:phyxminiproblems-problemphyxmini0334-sealedpressurizedtank, def:physics:phyx-mini-0334:phyxminiproblems-problemphyxmini0334-matchesproblemandfiguredata, def:physics:phyx-mini-0334:phyxminiproblems-problemphyxmini0334-usesstandardwaterandgravity, def:physics:phyx-mini-0334:phyxminiproblems-problemphyxmini0334-hascylindricalheadspacegeometry, def:physics:phyx-mini-0334:phyxminiproblems-problemphyxmini0334-satisfiesisothermalidealgaslaw, def:physics:phyx-mini-0334:phyxminiproblems-problemphyxmini0334-satisfieshydrostaticoutletlaw, def:physics:phyx-mini-0334:phyxminiproblems-problemphyxmini0334-flowstopsat, def:physics:phyx-mini-0334:phyxminiproblems-problemphyxmini0334-answerchoice, def:physics:phyx-mini-0334:phyxminiproblems-problemphyxmini0334-displayedheightmeters, def:physics:phyx-mini-0334:phyxminiproblems-problemphyxmini0334-recordeddatasetanswer, def:physics:phyx-mini-0334:phyxminiproblems-problemphyxmini0334-roundstodisplayedhundredthmeter}
The assigned autoformalize agent should translate this physics problem into Lean declarations in the covered file, with theorem and lemma proof bodies written as `by sorry`.
\end{theorem}
\begin{proof}
This is an autoformalization task, not a proof task. Produce statements that can later be proved by the physics prover without weakening the physical model.
\end{proof}

% --- Archon declaration topology begin ---
\section{Declaration topology}
\begin{definition}[Length Quantity]
  \label{def:physics:phyx-mini-0334:phyxminiproblems-problemphyxmini0334-lengthquantity}
  \lean{PhyXMiniProblems.ProblemPhyXMini0334.LengthQuantity}
  A nonnegative physical length
\end{definition}
\begin{proof}
  This is the declared model definition of Length Quantity.
\end{proof}
\begin{definition}[Area Quantity]
  \label{def:physics:phyx-mini-0334:phyxminiproblems-problemphyxmini0334-areaquantity}
  \lean{PhyXMiniProblems.ProblemPhyXMini0334.AreaQuantity}
  A nonnegative cross-sectional area
\end{definition}
\begin{proof}
  This is the declared model definition of Area Quantity.
\end{proof}
\begin{definition}[Volume Quantity]
  \label{def:physics:phyx-mini-0334:phyxminiproblems-problemphyxmini0334-volumequantity}
  \lean{PhyXMiniProblems.ProblemPhyXMini0334.VolumeQuantity}
  A nonnegative physical volume
\end{definition}
\begin{proof}
  This is the declared model definition of Volume Quantity.
\end{proof}
\begin{definition}[Mass Density Quantity]
  \label{def:physics:phyx-mini-0334:phyxminiproblems-problemphyxmini0334-massdensityquantity}
  \lean{PhyXMiniProblems.ProblemPhyXMini0334.MassDensityQuantity}
  A nonnegative mass density
\end{definition}
\begin{proof}
  This is the declared model definition of Mass Density Quantity.
\end{proof}
\begin{definition}[Acceleration Quantity]
  \label{def:physics:phyx-mini-0334:phyxminiproblems-problemphyxmini0334-accelerationquantity}
  \lean{PhyXMiniProblems.ProblemPhyXMini0334.AccelerationQuantity}
  A nonnegative acceleration magnitude
\end{definition}
\begin{proof}
  This is the declared model definition of Acceleration Quantity.
\end{proof}
\begin{definition}[Pressure Quantity]
  \label{def:physics:phyx-mini-0334:phyxminiproblems-problemphyxmini0334-pressurequantity}
  \lean{PhyXMiniProblems.ProblemPhyXMini0334.PressureQuantity}
  Absolute pressure, using Physlib's pressure dimension
\end{definition}
\begin{proof}
  This is the declared model definition of Pressure Quantity.
\end{proof}
\begin{definition}[nonnegative Quantity Readout]
  \label{def:physics:phyx-mini-0334:phyxminiproblems-problemphyxmini0334-nonnegativequantityreadout}
  \lean{PhyXMiniProblems.ProblemPhyXMini0334.nonnegativeQuantityReadout}
  Read a nonnegative dimensionful quantity in a coherent unit system
\end{definition}
\begin{proof}
  This is the declared model definition of nonnegative Quantity Readout.
\end{proof}
\begin{definition}[length In Meters]
  \label{def:physics:phyx-mini-0334:phyxminiproblems-problemphyxmini0334-lengthinmeters}
  \lean{PhyXMiniProblems.ProblemPhyXMini0334.lengthInMeters}
  \uses{def:physics:phyx-mini-0334:phyxminiproblems-problemphyxmini0334-lengthquantity, def:physics:phyx-mini-0334:phyxminiproblems-problemphyxmini0334-nonnegativequantityreadout}
  SI metre readout of a length
\end{definition}
\begin{proof}
  This is the declared model definition of length In Meters.
\end{proof}
\begin{definition}[area In Square Meters]
  \label{def:physics:phyx-mini-0334:phyxminiproblems-problemphyxmini0334-areainsquaremeters}
  \lean{PhyXMiniProblems.ProblemPhyXMini0334.areaInSquareMeters}
  \uses{def:physics:phyx-mini-0334:phyxminiproblems-problemphyxmini0334-areaquantity, def:physics:phyx-mini-0334:phyxminiproblems-problemphyxmini0334-nonnegativequantityreadout}
  SI square-metre readout of an area
\end{definition}
\begin{proof}
  This is the declared model definition of area In Square Meters.
\end{proof}
\begin{definition}[volume In Cubic Meters]
  \label{def:physics:phyx-mini-0334:phyxminiproblems-problemphyxmini0334-volumeincubicmeters}
  \lean{PhyXMiniProblems.ProblemPhyXMini0334.volumeInCubicMeters}
  \uses{def:physics:phyx-mini-0334:phyxminiproblems-problemphyxmini0334-volumequantity, def:physics:phyx-mini-0334:phyxminiproblems-problemphyxmini0334-nonnegativequantityreadout}
  SI cubic-metre readout of a volume
\end{definition}
\begin{proof}
  This is the declared model definition of volume In Cubic Meters.
\end{proof}
\begin{definition}[density In Kilograms Per Cubic Meter]
  \label{def:physics:phyx-mini-0334:phyxminiproblems-problemphyxmini0334-densityinkilogramspercubicmeter}
  \lean{PhyXMiniProblems.ProblemPhyXMini0334.densityInKilogramsPerCubicMeter}
  \uses{def:physics:phyx-mini-0334:phyxminiproblems-problemphyxmini0334-massdensityquantity, def:physics:phyx-mini-0334:phyxminiproblems-problemphyxmini0334-nonnegativequantityreadout}
  SI kilogram-per-cubic-metre readout of a mass density
\end{definition}
\begin{proof}
  This is the declared model definition of density In Kilograms Per Cubic Meter.
\end{proof}
\begin{definition}[acceleration In Meters Per Second Squared]
  \label{def:physics:phyx-mini-0334:phyxminiproblems-problemphyxmini0334-accelerationinmeterspersecondsquared}
  \lean{PhyXMiniProblems.ProblemPhyXMini0334.accelerationInMetersPerSecondSquared}
  \uses{def:physics:phyx-mini-0334:phyxminiproblems-problemphyxmini0334-accelerationquantity, def:physics:phyx-mini-0334:phyxminiproblems-problemphyxmini0334-nonnegativequantityreadout}
  SI metre-per-second-squared readout of an acceleration
\end{definition}
\begin{proof}
  This is the declared model definition of acceleration In Meters Per Second Squared.
\end{proof}
\begin{definition}[pressure In Pascals]
  \label{def:physics:phyx-mini-0334:phyxminiproblems-problemphyxmini0334-pressureinpascals}
  \lean{PhyXMiniProblems.ProblemPhyXMini0334.pressureInPascals}
  \uses{def:physics:phyx-mini-0334:phyxminiproblems-problemphyxmini0334-pressurequantity}
  SI pascal readout of an absolute pressure
\end{definition}
\begin{proof}
  This is the declared model definition of pressure In Pascals.
\end{proof}
\begin{definition}[Tank Shape]
  \label{def:physics:phyx-mini-0334:phyxminiproblems-problemphyxmini0334-tankshape}
  \lean{PhyXMiniProblems.ProblemPhyXMini0334.TankShape}
  Shape of the vessel shown in the supplied figure
\end{definition}
\begin{proof}
  This is the declared model definition of Tank Shape.
\end{proof}
\begin{definition}[Tank Top Boundary]
  \label{def:physics:phyx-mini-0334:phyxminiproblems-problemphyxmini0334-tanktopboundary}
  \lean{PhyXMiniProblems.ProblemPhyXMini0334.TankTopBoundary}
  Boundary condition at the top of the tank
\end{definition}
\begin{proof}
  This is the declared model definition of Tank Top Boundary.
\end{proof}
\begin{definition}[Tank Liquid]
  \label{def:physics:phyx-mini-0334:phyxminiproblems-problemphyxmini0334-tankliquid}
  \lean{PhyXMiniProblems.ProblemPhyXMini0334.TankLiquid}
  Liquid named in the problem statement
\end{definition}
\begin{proof}
  This is the declared model definition of Tank Liquid.
\end{proof}
\begin{definition}[Headspace Gas]
  \label{def:physics:phyx-mini-0334:phyxminiproblems-problemphyxmini0334-headspacegas}
  \lean{PhyXMiniProblems.ProblemPhyXMini0334.HeadspaceGas}
  Gas occupying the headspace above the water
\end{definition}
\begin{proof}
  This is the declared model definition of Headspace Gas.
\end{proof}
\begin{definition}[Outlet Environment]
  \label{def:physics:phyx-mini-0334:phyxminiproblems-problemphyxmini0334-outletenvironment}
  \lean{PhyXMiniProblems.ProblemPhyXMini0334.OutletEnvironment}
  Environment to which the side hose discharges
\end{definition}
\begin{proof}
  This is the declared model definition of Outlet Environment.
\end{proof}
\begin{definition}[Figure Quantity Label]
  \label{def:physics:phyx-mini-0334:phyxminiproblems-problemphyxmini0334-figurequantitylabel}
  \lean{PhyXMiniProblems.ProblemPhyXMini0334.FigureQuantityLabel}
  Quantity symbols printed in the primary figure
\end{definition}
\begin{proof}
  This is the declared model definition of Figure Quantity Label.
\end{proof}
\begin{definition}[Figure Quantity Role]
  \label{def:physics:phyx-mini-0334:phyxminiproblems-problemphyxmini0334-figurequantityrole}
  \lean{PhyXMiniProblems.ProblemPhyXMini0334.FigureQuantityRole}
  Physical role of each printed quantity symbol
\end{definition}
\begin{proof}
  This is the declared model definition of Figure Quantity Role.
\end{proof}
\begin{definition}[Figure Length Marker]
  \label{def:physics:phyx-mini-0334:phyxminiproblems-problemphyxmini0334-figurelengthmarker}
  \lean{PhyXMiniProblems.ProblemPhyXMini0334.FigureLengthMarker}
  The two explicit vertical length markers in the primary figure
\end{definition}
\begin{proof}
  This is the declared model definition of Figure Length Marker.
\end{proof}
\begin{definition}[Figure Location]
  \label{def:physics:phyx-mini-0334:phyxminiproblems-problemphyxmini0334-figurelocation}
  \lean{PhyXMiniProblems.ProblemPhyXMini0334.FigureLocation}
  Geometric location at which the visible hose is attached
\end{definition}
\begin{proof}
  This is the declared model definition of Figure Location.
\end{proof}
\begin{definition}[Tank Hose Figure]
  \label{def:physics:phyx-mini-0334:phyxminiproblems-problemphyxmini0334-tankhosefigure}
  \lean{PhyXMiniProblems.ProblemPhyXMini0334.TankHoseFigure}
  \uses{def:physics:phyx-mini-0334:phyxminiproblems-problemphyxmini0334-lengthquantity, def:physics:phyx-mini-0334:phyxminiproblems-problemphyxmini0334-tankshape, def:physics:phyx-mini-0334:phyxminiproblems-problemphyxmini0334-figurequantitylabel, def:physics:phyx-mini-0334:phyxminiproblems-problemphyxmini0334-figurequantityrole, def:physics:phyx-mini-0334:phyxminiproblems-problemphyxmini0334-figurelengthmarker, def:physics:phyx-mini-0334:phyxminiproblems-problemphyxmini0334-figurelocation}
  Structured transcription of the labels and geometry in image 334.png
\end{definition}
\begin{proof}
  This is the declared model definition of Tank Hose Figure.
\end{proof}
\begin{definition}[Sealed Pressurized Tank]
  \label{def:physics:phyx-mini-0334:phyxminiproblems-problemphyxmini0334-sealedpressurizedtank}
  \lean{PhyXMiniProblems.ProblemPhyXMini0334.SealedPressurizedTank}
  \uses{def:physics:phyx-mini-0334:phyxminiproblems-problemphyxmini0334-lengthquantity, def:physics:phyx-mini-0334:phyxminiproblems-problemphyxmini0334-areaquantity, def:physics:phyx-mini-0334:phyxminiproblems-problemphyxmini0334-volumequantity, def:physics:phyx-mini-0334:phyxminiproblems-problemphyxmini0334-massdensityquantity, def:physics:phyx-mini-0334:phyxminiproblems-problemphyxmini0334-accelerationquantity, def:physics:phyx-mini-0334:phyxminiproblems-problemphyxmini0334-pressurequantity, def:physics:phyx-mini-0334:phyxminiproblems-problemphyxmini0334-tanktopboundary, def:physics:phyx-mini-0334:phyxminiproblems-problemphyxmini0334-tankliquid, def:physics:phyx-mini-0334:phyxminiproblems-problemphyxmini0334-headspacegas, def:physics:phyx-mini-0334:phyxminiproblems-problemphyxmini0334-outletenvironment, def:physics:phyx-mini-0334:phyxminiproblems-problemphyxmini0334-tankhosefigure}
  The physical apparatus and state-dependent observables. Pressure, headspace volume, absolute temperature, and liquid pressure at the inside of the outlet are independent fields here. Their physical relations are supplied by the governing-law predicates below; in particular, the stopped height and the recorded answer are not definitions in this setup
\end{definition}
\begin{proof}
  This is the declared model definition of Sealed Pressurized Tank.
\end{proof}
\begin{definition}[Is Admissible Water Height]
  \label{def:physics:phyx-mini-0334:phyxminiproblems-problemphyxmini0334-isadmissiblewaterheight}
  \lean{PhyXMiniProblems.ProblemPhyXMini0334.IsAdmissibleWaterHeight}
  \uses{def:physics:phyx-mini-0334:phyxminiproblems-problemphyxmini0334-lengthquantity, def:physics:phyx-mini-0334:phyxminiproblems-problemphyxmini0334-lengthinmeters, def:physics:phyx-mini-0334:phyxminiproblems-problemphyxmini0334-sealedpressurizedtank}
  A water height lies inside the cylindrical vessel
\end{definition}
\begin{proof}
  This is the declared model definition of Is Admissible Water Height.
\end{proof}
\begin{definition}[Occurs Along Outflow From Initial State]
  \label{def:physics:phyx-mini-0334:phyxminiproblems-problemphyxmini0334-occursalongoutflowfrominitialstate}
  \lean{PhyXMiniProblems.ProblemPhyXMini0334.OccursAlongOutflowFromInitialState}
  \uses{def:physics:phyx-mini-0334:phyxminiproblems-problemphyxmini0334-lengthquantity, def:physics:phyx-mini-0334:phyxminiproblems-problemphyxmini0334-lengthinmeters, def:physics:phyx-mini-0334:phyxminiproblems-problemphyxmini0334-sealedpressurizedtank, def:physics:phyx-mini-0334:phyxminiproblems-problemphyxmini0334-isadmissiblewaterheight}
  A height is on the decreasing-height branch reached after discharge begins: the outlet is still covered and the height has not increased above its initial value
\end{definition}
\begin{proof}
  This is the declared model definition of Occurs Along Outflow From Initial State.
\end{proof}
\begin{definition}[Matches Problem And Figure Data]
  \label{def:physics:phyx-mini-0334:phyxminiproblems-problemphyxmini0334-matchesproblemandfiguredata}
  \lean{PhyXMiniProblems.ProblemPhyXMini0334.MatchesProblemAndFigureData}
  \uses{def:physics:phyx-mini-0334:phyxminiproblems-problemphyxmini0334-lengthinmeters, def:physics:phyx-mini-0334:phyxminiproblems-problemphyxmini0334-pressureinpascals, def:physics:phyx-mini-0334:phyxminiproblems-problemphyxmini0334-sealedpressurizedtank}
  The stated calibration and the readouts visible in the primary figure. The pressure values are absolute. The figure's p labels the compressed-air pressure at the water surface, while h spans from the tank base to that surface. No stopped height or answer choice is included here
\end{definition}
\begin{proof}
  This is the declared model definition of Matches Problem And Figure Data.
\end{proof}
\begin{definition}[Uses Standard Water And Gravity]
  \label{def:physics:phyx-mini-0334:phyxminiproblems-problemphyxmini0334-usesstandardwaterandgravity}
  \lean{PhyXMiniProblems.ProblemPhyXMini0334.UsesStandardWaterAndGravity}
  \uses{def:physics:phyx-mini-0334:phyxminiproblems-problemphyxmini0334-densityinkilogramspercubicmeter, def:physics:phyx-mini-0334:phyxminiproblems-problemphyxmini0334-accelerationinmeterspersecondsquared, def:physics:phyx-mini-0334:phyxminiproblems-problemphyxmini0334-sealedpressurizedtank}
  Standard near-surface numerical model for water on Earth, in coherent SI units. These implicit textbook constants are separated from the values explicitly printed in the problem
\end{definition}
\begin{proof}
  This is the declared model definition of Uses Standard Water And Gravity.
\end{proof}
\begin{definition}[Has Cylindrical Headspace Geometry]
  \label{def:physics:phyx-mini-0334:phyxminiproblems-problemphyxmini0334-hascylindricalheadspacegeometry}
  \lean{PhyXMiniProblems.ProblemPhyXMini0334.HasCylindricalHeadspaceGeometry}
  \uses{def:physics:phyx-mini-0334:phyxminiproblems-problemphyxmini0334-lengthquantity, def:physics:phyx-mini-0334:phyxminiproblems-problemphyxmini0334-lengthinmeters, def:physics:phyx-mini-0334:phyxminiproblems-problemphyxmini0334-areainsquaremeters, def:physics:phyx-mini-0334:phyxminiproblems-problemphyxmini0334-volumeincubicmeters, def:physics:phyx-mini-0334:phyxminiproblems-problemphyxmini0334-sealedpressurizedtank, def:physics:phyx-mini-0334:phyxminiproblems-problemphyxmini0334-isadmissiblewaterheight}
  For the cylindrical tank, headspace volume equals cross-sectional area times the vertical distance from the water surface to the sealed top
\end{definition}
\begin{proof}
  This is the declared model definition of Has Cylindrical Headspace Geometry.
\end{proof}
\begin{definition}[Satisfies Isothermal Ideal Gas Law]
  \label{def:physics:phyx-mini-0334:phyxminiproblems-problemphyxmini0334-satisfiesisothermalidealgaslaw}
  \lean{PhyXMiniProblems.ProblemPhyXMini0334.SatisfiesIsothermalIdealGasLaw}
  \uses{def:physics:phyx-mini-0334:phyxminiproblems-problemphyxmini0334-lengthquantity, def:physics:phyx-mini-0334:phyxminiproblems-problemphyxmini0334-volumeincubicmeters, def:physics:phyx-mini-0334:phyxminiproblems-problemphyxmini0334-pressureinpascals, def:physics:phyx-mini-0334:phyxminiproblems-problemphyxmini0334-sealedpressurizedtank, def:physics:phyx-mini-0334:phyxminiproblems-problemphyxmini0334-isadmissiblewaterheight}
  The sealed headspace follows an isothermal ideal-gas process. The first field states constant absolute temperature; the second is Boyle's law p V = constant for any two admissible water heights
\end{definition}
\begin{proof}
  This is the declared model definition of Satisfies Isothermal Ideal Gas Law.
\end{proof}
\begin{definition}[Satisfies Hydrostatic Outlet Law]
  \label{def:physics:phyx-mini-0334:phyxminiproblems-problemphyxmini0334-satisfieshydrostaticoutletlaw}
  \lean{PhyXMiniProblems.ProblemPhyXMini0334.SatisfiesHydrostaticOutletLaw}
  \uses{def:physics:phyx-mini-0334:phyxminiproblems-problemphyxmini0334-lengthquantity, def:physics:phyx-mini-0334:phyxminiproblems-problemphyxmini0334-lengthinmeters, def:physics:phyx-mini-0334:phyxminiproblems-problemphyxmini0334-densityinkilogramspercubicmeter, def:physics:phyx-mini-0334:phyxminiproblems-problemphyxmini0334-accelerationinmeterspersecondsquared, def:physics:phyx-mini-0334:phyxminiproblems-problemphyxmini0334-pressureinpascals, def:physics:phyx-mini-0334:phyxminiproblems-problemphyxmini0334-sealedpressurizedtank, def:physics:phyx-mini-0334:phyxminiproblems-problemphyxmini0334-isadmissiblewaterheight}
  Hydrostatic pressure at the inside of the side outlet equals the air pressure at the water surface plus rho * g times the water depth above the outlet
\end{definition}
\begin{proof}
  This is the declared model definition of Satisfies Hydrostatic Outlet Law.
\end{proof}
\begin{definition}[Flow Stops At]
  \label{def:physics:phyx-mini-0334:phyxminiproblems-problemphyxmini0334-flowstopsat}
  \lean{PhyXMiniProblems.ProblemPhyXMini0334.FlowStopsAt}
  \uses{def:physics:phyx-mini-0334:phyxminiproblems-problemphyxmini0334-lengthquantity, def:physics:phyx-mini-0334:phyxminiproblems-problemphyxmini0334-pressureinpascals, def:physics:phyx-mini-0334:phyxminiproblems-problemphyxmini0334-sealedpressurizedtank, def:physics:phyx-mini-0334:phyxminiproblems-problemphyxmini0334-occursalongoutflowfrominitialstate}
  Flow stops on the physical decreasing-height branch when the static liquid pressure just inside the hose equals the atmospheric pressure outside it
\end{definition}
\begin{proof}
  This is the declared model definition of Flow Stops At.
\end{proof}
\begin{lemma}[air pressure times headspace height eq initial]
  \label{lem:physics:phyx-mini-0334:phyxminiproblems-problemphyxmini0334-air-pressure-times-headspace-height-eq-initial}
  \lean{PhyXMiniProblems.ProblemPhyXMini0334.air_pressure_times_headspace_height_eq_initial}
  \uses{def:physics:phyx-mini-0334:phyxminiproblems-problemphyxmini0334-lengthquantity, def:physics:phyx-mini-0334:phyxminiproblems-problemphyxmini0334-lengthinmeters, def:physics:phyx-mini-0334:phyxminiproblems-problemphyxmini0334-pressureinpascals, def:physics:phyx-mini-0334:phyxminiproblems-problemphyxmini0334-sealedpressurizedtank, def:physics:phyx-mini-0334:phyxminiproblems-problemphyxmini0334-isadmissiblewaterheight, def:physics:phyx-mini-0334:phyxminiproblems-problemphyxmini0334-hascylindricalheadspacegeometry, def:physics:phyx-mini-0334:phyxminiproblems-problemphyxmini0334-satisfiesisothermalidealgaslaw}
  Boyle's law and cylindrical geometry give the pressure-height relation
\end{lemma}
\begin{proof}
  The conclusion follows from the governing relations and figure readouts named in the statement of air pressure times headspace height eq initial.
\end{proof}
\begin{lemma}[stopped height pressure balance]
  \label{lem:physics:phyx-mini-0334:phyxminiproblems-problemphyxmini0334-stopped-height-pressure-balance}
  \lean{PhyXMiniProblems.ProblemPhyXMini0334.stopped_height_pressure_balance}
  \uses{def:physics:phyx-mini-0334:phyxminiproblems-problemphyxmini0334-lengthquantity, def:physics:phyx-mini-0334:phyxminiproblems-problemphyxmini0334-lengthinmeters, def:physics:phyx-mini-0334:phyxminiproblems-problemphyxmini0334-densityinkilogramspercubicmeter, def:physics:phyx-mini-0334:phyxminiproblems-problemphyxmini0334-accelerationinmeterspersecondsquared, def:physics:phyx-mini-0334:phyxminiproblems-problemphyxmini0334-pressureinpascals, def:physics:phyx-mini-0334:phyxminiproblems-problemphyxmini0334-sealedpressurizedtank, def:physics:phyx-mini-0334:phyxminiproblems-problemphyxmini0334-satisfieshydrostaticoutletlaw, def:physics:phyx-mini-0334:phyxminiproblems-problemphyxmini0334-flowstopsat}
  At a stopped height, hydrostatic outlet pressure balances atmosphere
\end{lemma}
\begin{proof}
  The conclusion follows from the governing relations and figure readouts named in the statement of stopped height pressure balance.
\end{proof}
\begin{lemma}[stopped height satisfies quadratic]
  \label{lem:physics:phyx-mini-0334:phyxminiproblems-problemphyxmini0334-stopped-height-satisfies-quadratic}
  \lean{PhyXMiniProblems.ProblemPhyXMini0334.stopped_height_satisfies_quadratic}
  \uses{def:physics:phyx-mini-0334:phyxminiproblems-problemphyxmini0334-lengthquantity, def:physics:phyx-mini-0334:phyxminiproblems-problemphyxmini0334-lengthinmeters, def:physics:phyx-mini-0334:phyxminiproblems-problemphyxmini0334-sealedpressurizedtank, def:physics:phyx-mini-0334:phyxminiproblems-problemphyxmini0334-matchesproblemandfiguredata, def:physics:phyx-mini-0334:phyxminiproblems-problemphyxmini0334-usesstandardwaterandgravity, def:physics:phyx-mini-0334:phyxminiproblems-problemphyxmini0334-hascylindricalheadspacegeometry, def:physics:phyx-mini-0334:phyxminiproblems-problemphyxmini0334-satisfiesisothermalidealgaslaw, def:physics:phyx-mini-0334:phyxminiproblems-problemphyxmini0334-satisfieshydrostaticoutletlaw, def:physics:phyx-mini-0334:phyxminiproblems-problemphyxmini0334-flowstopsat}
  After inserting the printed data and standard water/gravity values, the physical stopped height satisfies the resulting quadratic equation
\end{lemma}
\begin{proof}
  The conclusion follows from the governing relations and figure readouts named in the statement of stopped height satisfies quadratic.
\end{proof}
\begin{definition}[Answer Choice]
  \label{def:physics:phyx-mini-0334:phyxminiproblems-problemphyxmini0334-answerchoice}
  \lean{PhyXMiniProblems.ProblemPhyXMini0334.AnswerChoice}
  The four height choices printed in the problem
\end{definition}
\begin{proof}
  This is the declared model definition of Answer Choice.
\end{proof}
\begin{definition}[displayed Height Meters]
  \label{def:physics:phyx-mini-0334:phyxminiproblems-problemphyxmini0334-displayedheightmeters}
  \lean{PhyXMiniProblems.ProblemPhyXMini0334.displayedHeightMeters}
  \uses{def:physics:phyx-mini-0334:phyxminiproblems-problemphyxmini0334-answerchoice}
  Metre readout displayed beside an answer choice
\end{definition}
\begin{proof}
  This is the declared model definition of displayed Height Meters.
\end{proof}
\begin{definition}[recorded Dataset Answer]
  \label{def:physics:phyx-mini-0334:phyxminiproblems-problemphyxmini0334-recordeddatasetanswer}
  \lean{PhyXMiniProblems.ProblemPhyXMini0334.recordedDatasetAnswer}
  \uses{def:physics:phyx-mini-0334:phyxminiproblems-problemphyxmini0334-answerchoice}
  The answer label recorded in the dataset metadata
\end{definition}
\begin{proof}
  This is the declared model definition of recorded Dataset Answer.
\end{proof}
\begin{definition}[Rounds To Displayed Hundredth Meter]
  \label{def:physics:phyx-mini-0334:phyxminiproblems-problemphyxmini0334-roundstodisplayedhundredthmeter}
  \lean{PhyXMiniProblems.ProblemPhyXMini0334.RoundsToDisplayedHundredthMeter}
  A physical height rounds to the displayed value at hundredth-metre precision
\end{definition}
\begin{proof}
  This is the declared model definition of Rounds To Displayed Hundredth Meter.
\end{proof}
% --- Archon declaration topology end ---
% --- Archon physics formalization source end ---
